\chapter{Introduction}\label{ch:intro}
\section{Research context}
Networks describe interactions whose consequences extend beyond individual records. Selecting a few influential starting nodes is useful when resources limit how broadly an intervention can be introduced. Influence maximization formalizes this decision as the choice of a seed set under a diffusion model. The classical formulation connects this problem to submodular optimization and gives a principled basis for greedy approximation under suitable assumptions \cite{kempe2003}. Such guarantees concern the stated diffusion objective; they do not establish that a learned score or a different implementation has the same approximation property.

Graph neural networks provide another route: learn a structural scoring function and rank candidate nodes. Their attraction is that a parameterized computation can transfer structural patterns from a training graph to a held-out graph. The present study starts from the existing PrivIM reproduction rather than introducing a new GNN architecture. It retains the released-code GRAT3 model, structural input features, original loss, subgraph sampling, and fixed evaluation rules. This makes the privacy mechanism and its empirical consequences the central objects of investigation.

Graph learning raises a difficult privacy question because one person's presence can affect many computational contributions. Removing a node removes its incident edges and may change walks, sampled subgraphs, message passing, and optimization. Clipping a gradient for each nominal training item is insufficient unless the number of items influenced by a node is also controlled. Privacy therefore depends on the relation between graph adjacency and the entire training query, rather than on a noise parameter considered in isolation.

\section{Heterogeneous protection requirements}
A single privacy budget is convenient but may be conservative when participants have different requirements. In the setting studied here, each fixed public identity has a target drawn from $\{1,2,4\}$, with a common target $\delta=10^{-5}$. Smaller epsilon represents a more demanding permitted bound. Uniform Node-DP calibrated to epsilon 1 provides a strict policy control compatible with all three target classes. It also leaves unused allowance for identities assigned less restrictive targets.

Personalization offers a possible tradeoff. A node with a stricter requirement may participate less often in a mechanism whose base privacy parameter is larger. This can permit lower Gaussian noise for training while satisfying heterogeneous caps. The potential gain has a cost: omitting nodes also omits their incident structure, potentially reducing the information available to learn useful rankings. Neither a privacy amplification formula nor the existence of larger budgets determines which effect dominates Coverage.

This distinction motivates the research question: \emph{under heterogeneous node requirements, can a personalized node-private GNN training mechanism improve Coverage over a uniform mechanism calibrated to the strictest policy?} The question is empirical as well as mathematical. A mechanism may satisfy a valid conditional privacy guarantee and still fail to deliver a reliable utility advantage.

\section{Objectives and scope}
The first objective is to formulate a node-level personalized privacy problem with explicit identity, adjacency, and release assumptions. The second is to describe and analyze the PDP-PrivIM mechanism: personalized participation, stable overlapping slots, double clipping, full-coordinate Gaussian optimization, and a fixed last-layer projection. The third is to compare Coverage and observed variability across five datasets and four configurations, reporting three runs per dataset/configuration cell, or 60 runs in total. Historical experiments provide supplementary evidence about training variability and are distinguished from the current three-run comparison.

The four configurations have different roles. PDP-PrivIM is the personalized mechanism. Uniform epsilon 1 is the primary strict-policy control. Uniform epsilon 4 is a relaxed-policy reference; it is not certified by its stated bound to satisfy the original target 1 and target 2 caps. Non-Private is a performance reference without a DP guarantee. The comparison consequently does not assume an ordering among all four methods, nor interpret every difference as the isolated effect of personalization.

The scope is deliberately fixed. Training and test graphs use the existing prepared partitions. The GNN architecture, original objective, IC implementation, SCS and BES are preserved. The study evaluates Top-50 selection with the original one-step outgoing IC rule at probability one. This differs from unrestricted multi-round diffusion and limits the interpretation of the results. No claim is made that the five graphs represent a random sample of all social networks.

The privacy theorem concerns a fixed prepared-domain training mechanism and a specified model output. Raw preprocessing, publishing an identity universe, private test evaluation, resource side channels, finite-precision implementation, and joint release of many research models require separate analysis. An experimental result is not automatically a private public release.

\section{Evidence and principal findings}
The empirical chapter reports five datasets and four configurations, with three runs per cell and 60 runs in total. Every cell gives Coverage mean and sample SD. PDP-PrivIM has a higher observed mean than Uniform1 on all five datasets, remains below Non-Private on all five, and exceeds Uniform4 only on LastFM. Its sample SD is lower than Uniform1 on two datasets and higher on three. These are descriptive results from existing development experiments; their source protocol is documented in Appendix~\ref{app:development}.

The main comparison addresses the observed utility tradeoff under heterogeneous privacy requirements. Historical Gaussian-path diagnostics, static update scales, and complete five-seed outcomes separately inform the discussion of training variability. Their observations do not enlarge the current sample size or establish a universal performance ordering.

\section{Candidate contributions and attribution}
The dissertation presents three evidence-supported candidate contributions. First, it formulates and implements a specific combination of established tools for heterogeneous node-private GNN influence maximization. Second, it organizes a conditional privacy argument connecting stable slot locality, sensitivity, adaptive Gaussian composition, and personalized participation. Third, it characterizes the utility, uncertainty, and low-tail behavior of the frozen mechanism using three runs per dataset/configuration cell, with historical variability evidence documented separately.

I carried out the design and implementation of the personalized mechanism, developed its privacy analysis, and conducted and analyzed the experiments. This work builds on the existing PrivIM reproduction, which provides the GRAT3 architecture, training objective, and influence-maximization sampling components described in Chapter~\ref{ch:related}. The privacy analysis also uses established tools for gradient clipping, Gaussian randomization, RDP, and privacy amplification.

The empirical evidence supports an observed mean-Coverage advantage over the strict uniform control in the reported three-run comparison, but does not establish that personalization improves expected utility or stability in general. The historical evidence also records substantial private-training variability. Both findings matter when interpreting a later independent evaluation.

\section{Dissertation organization}
Chapter~\ref{ch:related} reviews influence maximization, private learning, and personalized graph privacy. Chapter~\ref{ch:problem} defines the graph problem, threat model, policies, and estimand. Chapter~\ref{ch:method} describes the frozen mechanism and conditional proof. Chapter~\ref{ch:design} sets out the experimental and statistical methodology. Chapter~\ref{ch:results} reports the three-run comparison and distinguishes historical five-seed evidence and controlled diagnostics. Chapter~\ref{ch:conclusion} concludes with limitations and future research directions. Appendices give reproducibility and historical development details with current individual outcomes and historical records traceable.
